\chapter{Installation and Dependencies}

\section{The Package}

Originally, the \ghats{} package could be obtained from INAF-OAB and was not
distributed broadly.  This was because Tomaso did not have the resources to
support the software for a large community.  From version 3.3.0, the code is
available from the GitHub repository so that the community can use it, test it,
and contribute fixes or contributed routines on a development branch.

\section{Unpacking}

You can unpack or clone the package in your preferred directory.  The
IDL/GDL part of \ghats{} is distributed as source code, so the tree is made of
ASCII routines rather than compiled binaries.  A typical checkout contains the
core directories

\begin{itemize}
  \item \param{FFT}: routines used in FFT and PDS production;
  \item \param{ANALYSIS}: routines for interactive analysis of GHATS products;
  \item \param{CROSS_ANALYSIS}: routines dedicated to cross spectra, lags, and
        coherence;
  \item \param{CONTRIBUTED}: contributed routines;
  \item \param{PYTHON}: Python-based GHATS tools, including the command-line
        helper directory;
  \item \param{DOCS}: documentation and manual sources.
\end{itemize}

The tree also contains mission-specific directories, for example
\param{NICER}, \param{NuSTAR}, \param{Swift}, \param{XMM}, \param{IXPE},
\param{ASTROSAT}, \param{Huiyan}, and \param{Einstein_Probe}.  These contain
the routines that know how to read the native event products of the corresponding
instruments and produce GHATS-format PDS or FFT files.

If your local installation includes launcher scripts such as \param{ghats} for
IDL or \param{gghats} for GDL, put the one you use in a directory that is in your
shell \param{PATH}, and make it executable.  Otherwise, the same setup can be
done directly in your shell by defining the environment variables described
below and then starting IDL or GDL.

\section{Software Dependencies}

You need the following software and libraries installed:

\begin{enumerate}
  \item An updated and running version of IDL or GDL.  If you use a precompiled
        GDL for your architecture, you may still need the source distribution,
        because many procedures are external.  GDL can be installed following
        the instructions at
        \url{https://github.com/gnudatalanguage/gdl/wiki/GDL-on-Linux}.
  \item An updated version of the IDL Astronomy User's Library:
        \url{https://asd.gsfc.nasa.gov/archive/idlastro/}.
  \item An updated version of Craig Markwardt's IDL library:
        \url{https://pages.physics.wisc.edu/~craigm/idl/idl.html}.
\end{enumerate}

Not every routine in those libraries is required by every GHATS command, but
they provide many utilities that the package expects to have available.

\section{Environment Setup}

To run \ghats{}, your shell must know where the repository is, where IDL or GDL
is installed, and where the required IDL libraries are.  The key variables are
\param{MU_PATH}, \param{IDL_STARTUP}, \param{IDL_DIR}, and \param{IDL_PATH}.

For a Bourne-like shell such as \param{bash}, a typical IDL setup is

\begin{ghatsconsole}
export MU_PATH=/path/to/GHATS
export IDL_STARTUP=$MU_PATH/ANALYSIS/ghats.pro
export IDL_DIR=/path/to/idl
export IDL_PATH=+$MU_PATH:+/path/to/astron.dir:+/path/to/markwardt:+$IDL_DIR/lib
export PATH=$PATH:$IDL_DIR/bin
export PATH=$PATH:$MU_PATH/PYTHON
\end{ghatsconsole}

Replace the paths with the locations on your machine.  The leading \param{+} in
\param{IDL_PATH} tells IDL/GDL to search a directory and its subdirectories.

For a C-shell-like shell such as \param{tcsh}, the same setup uses
\param{setenv}:

\begin{ghatsconsole}
setenv MU_PATH /path/to/GHATS
setenv IDL_STARTUP $MU_PATH/ANALYSIS/ghats.pro
setenv IDL_DIR /path/to/idl
setenv IDL_PATH +$MU_PATH:+/path/to/astron.dir:+/path/to/markwardt:+$IDL_DIR/lib
setenv PATH ${PATH}:$IDL_DIR/bin:$MU_PATH/PYTHON
\end{ghatsconsole}

These commands can be put in your shell startup file if you want them to be
available in every terminal session.  For example, \param{bash} users typically
use \param{.bashrc} or \param{.bash_profile}, while \param{tcsh} users commonly
use \param{.login}.

\section{Running the Program}

After the environment is set, start the appropriate interpreter or launcher.  If
you use a launcher script, the commands are typically

\begin{ghatsconsole}
$ ghats
\end{ghatsconsole}

for IDL, or

\begin{ghatsconsole}
$ gghats
\end{ghatsconsole}

for GDL.  If you start IDL directly, the startup file should load the GHATS
environment and change the prompt to

\begin{ghatsconsole}
Ghats>
\end{ghatsconsole}

At that point the GHATS routines can be called directly.  The following chapters
give examples of producing PDS and FFT files, reading them back, plotting timing
products, and using the cross-spectral tools.
